% Exact two-outcome L2 geometry: probabilities (1/2,1/2), trivial information.
\begin{tikzpicture}[
 x=1mm,y=1mm,
 every node/.style={font=\figurefont\fontsize{8}{10}\selectfont,text=FigureInk,inner sep=1pt},
 axis/.style={draw=FigureStroke,line width=.6pt,->},
 guide/.style={draw=FigureGrid,line width=.5pt,dashed}
]
 \path[use as bounding box] (0,0) rectangle (92,67);
 \draw[axis] (10,10) -- (82,10) node[below left=2pt] {第一结果的取值};
 \draw[axis] (10,10) -- (10,61);
 \node[rotate=90] at (2.5,38) {第二结果的取值};
 \draw[draw=FigureBlue,line width=1.1pt] (10,10) -- (57,57);
 \node[text=FigureBlue,align=left,anchor=west] at (60,55) {允许的预测\\\(U=(c,c)\)};
 \coordinate (X) at (10,50);
 \coordinate (M) at (30,30);
 \coordinate (U) at (50,50);
 \draw[draw=FigureGreen,line width=1.1pt] (X) -- (M);
 \draw[draw=FigureBlue,line width=1.1pt] (M) -- (U);
 \draw[draw=FigureMuted,line width=.75pt,dashed] (X) -- (U);
 \draw[draw=FigureStroke,line width=.55pt] (27.8,32.2) -- (30,34.4) -- (32.2,32.2);
 \fill[FigureInk] (X) circle (1mm);
 \fill[FigureGreen] (M) circle (1mm);
 \fill[FigureBlue] (U) circle (1mm);
 \node[above right=2pt] at (X) {\(X=(0,2)\)};
 \node[below right=3pt] at (M) {\(m=(1,1)\)};
 \node[below right=3pt] at (U) {\(U=(2,2)\)};
 \node[text=FigureGreen,rotate=-45] at (16,34) {\(X-m\)};
 \node[text=FigureBlue,rotate=45] at (45,38) {\(m-U\)};
 \node[text=FigureMuted] at (30,56) {\(X-U\)};
 \node[below=2pt] at (10,10) {\(0\)};
 \draw[guide] (30,10) -- (30,29);
 \node[below=2pt] at (30,10) {\(1\)};
 \draw[guide] (50,10) -- (50,49);
 \node[below=2pt] at (50,10) {\(2\)};
\end{tikzpicture}
